\documentclass[11pt]{article}
\newcommand{\ListaTitulo}{Lista 12 --- Fatoriais A×B e A×B×C}
% Preâmbulo compartilhado — MATD48, Listas de Exercícios 2026
% Usado por Lista{NN}.tex e Gabarito{NN}.tex via % Preâmbulo compartilhado — MATD48, Listas de Exercícios 2026
% Usado por Lista{NN}.tex e Gabarito{NN}.tex via % Preâmbulo compartilhado — MATD48, Listas de Exercícios 2026
% Usado por Lista{NN}.tex e Gabarito{NN}.tex via \input{preamble}

\usepackage[utf8]{inputenc}
\usepackage[T1]{fontenc}
\usepackage[brazil]{babel}
\usepackage{fullpage}
\usepackage{amsmath,amssymb}
\usepackage{enumitem}
\usepackage{xcolor}
\usepackage{fancyhdr}
\usepackage{titlesec}
\usepackage{listings}
\usepackage{hyperref}
\usepackage{booktabs}
\usepackage{graphicx}

\definecolor{matd48azul}{HTML}{0B4F6C}
\definecolor{matd48verde}{HTML}{2E8B57}

\titleformat{\section}{\normalfont\Large\bfseries\color{matd48azul}}{\thesection}{1em}{}
\titleformat{\subsection}{\normalfont\large\bfseries\color{matd48azul}}{\thesubsection}{1em}{}

\providecommand{\ListaTitulo}{Lista de Exercícios}

\setlength{\headheight}{14pt}
\pagestyle{fancy}
\fancyhf{}
\lhead{\small MATD48 --- Planejamento de Experimentos A}
\rhead{\small \ListaTitulo}
\cfoot{\thepage}
\renewcommand{\headrulewidth}{0.4pt}

\lstset{
  basicstyle=\ttfamily\small,
  breaklines=true,
  frame=single,
  columns=fullflexible,
  backgroundcolor=\color{gray!8},
  commentstyle=\color{matd48verde},
  keywordstyle=\color{matd48azul}\bfseries,
  language=R,
  extendedchars=true,
  literate=%
    {á}{{\'a}}1 {à}{{\`a}}1 {â}{{\^a}}1 {ã}{{\~a}}1
    {é}{{\'e}}1 {ê}{{\^e}}1 {í}{{\'i}}1 {ó}{{\'o}}1
    {ô}{{\^o}}1 {õ}{{\~o}}1 {ú}{{\'u}}1 {ç}{{\c c}}1
    {Á}{{\'A}}1 {À}{{\`A}}1 {Â}{{\^A}}1 {Ã}{{\~A}}1
    {É}{{\'E}}1 {Ê}{{\^E}}1 {Í}{{\'I}}1 {Ó}{{\'O}}1
    {Ô}{{\^O}}1 {Õ}{{\~O}}1 {Ú}{{\'U}}1 {Ç}{{\c C}}1
}

\hypersetup{colorlinks=true, linkcolor=matd48azul, urlcolor=matd48azul, citecolor=matd48azul}

\newcommand{\cabecalho}[2]{%
  \begin{center}
    {\Large\bfseries\color{matd48azul} #1}\\[0.3em]
    {\large #2}\\[0.3em]
    \rule{\linewidth}{0.6pt}
  \end{center}
  \vspace{0.5em}
}

\newenvironment{questao}[1]{\vspace{1em}\noindent\textbf{Questão #1.}\hspace{0.4em}}{\par}
\newenvironment{solucao}{\vspace{0.3em}\noindent\textit{Solução.}\hspace{0.4em}\color{black!85}}{\par\vspace{0.5em}}


\usepackage[utf8]{inputenc}
\usepackage[T1]{fontenc}
\usepackage[brazil]{babel}
\usepackage{fullpage}
\usepackage{amsmath,amssymb}
\usepackage{enumitem}
\usepackage{xcolor}
\usepackage{fancyhdr}
\usepackage{titlesec}
\usepackage{listings}
\usepackage{hyperref}
\usepackage{booktabs}
\usepackage{graphicx}

\definecolor{matd48azul}{HTML}{0B4F6C}
\definecolor{matd48verde}{HTML}{2E8B57}

\titleformat{\section}{\normalfont\Large\bfseries\color{matd48azul}}{\thesection}{1em}{}
\titleformat{\subsection}{\normalfont\large\bfseries\color{matd48azul}}{\thesubsection}{1em}{}

\providecommand{\ListaTitulo}{Lista de Exercícios}

\setlength{\headheight}{14pt}
\pagestyle{fancy}
\fancyhf{}
\lhead{\small MATD48 --- Planejamento de Experimentos A}
\rhead{\small \ListaTitulo}
\cfoot{\thepage}
\renewcommand{\headrulewidth}{0.4pt}

\lstset{
  basicstyle=\ttfamily\small,
  breaklines=true,
  frame=single,
  columns=fullflexible,
  backgroundcolor=\color{gray!8},
  commentstyle=\color{matd48verde},
  keywordstyle=\color{matd48azul}\bfseries,
  language=R,
  extendedchars=true,
  literate=%
    {á}{{\'a}}1 {à}{{\`a}}1 {â}{{\^a}}1 {ã}{{\~a}}1
    {é}{{\'e}}1 {ê}{{\^e}}1 {í}{{\'i}}1 {ó}{{\'o}}1
    {ô}{{\^o}}1 {õ}{{\~o}}1 {ú}{{\'u}}1 {ç}{{\c c}}1
    {Á}{{\'A}}1 {À}{{\`A}}1 {Â}{{\^A}}1 {Ã}{{\~A}}1
    {É}{{\'E}}1 {Ê}{{\^E}}1 {Í}{{\'I}}1 {Ó}{{\'O}}1
    {Ô}{{\^O}}1 {Õ}{{\~O}}1 {Ú}{{\'U}}1 {Ç}{{\c C}}1
}

\hypersetup{colorlinks=true, linkcolor=matd48azul, urlcolor=matd48azul, citecolor=matd48azul}

\newcommand{\cabecalho}[2]{%
  \begin{center}
    {\Large\bfseries\color{matd48azul} #1}\\[0.3em]
    {\large #2}\\[0.3em]
    \rule{\linewidth}{0.6pt}
  \end{center}
  \vspace{0.5em}
}

\newenvironment{questao}[1]{\vspace{1em}\noindent\textbf{Questão #1.}\hspace{0.4em}}{\par}
\newenvironment{solucao}{\vspace{0.3em}\noindent\textit{Solução.}\hspace{0.4em}\color{black!85}}{\par\vspace{0.5em}}


\usepackage[utf8]{inputenc}
\usepackage[T1]{fontenc}
\usepackage[brazil]{babel}
\usepackage{fullpage}
\usepackage{amsmath,amssymb}
\usepackage{enumitem}
\usepackage{xcolor}
\usepackage{fancyhdr}
\usepackage{titlesec}
\usepackage{listings}
\usepackage{hyperref}
\usepackage{booktabs}
\usepackage{graphicx}

\definecolor{matd48azul}{HTML}{0B4F6C}
\definecolor{matd48verde}{HTML}{2E8B57}

\titleformat{\section}{\normalfont\Large\bfseries\color{matd48azul}}{\thesection}{1em}{}
\titleformat{\subsection}{\normalfont\large\bfseries\color{matd48azul}}{\thesubsection}{1em}{}

\providecommand{\ListaTitulo}{Lista de Exercícios}

\setlength{\headheight}{14pt}
\pagestyle{fancy}
\fancyhf{}
\lhead{\small MATD48 --- Planejamento de Experimentos A}
\rhead{\small \ListaTitulo}
\cfoot{\thepage}
\renewcommand{\headrulewidth}{0.4pt}

\lstset{
  basicstyle=\ttfamily\small,
  breaklines=true,
  frame=single,
  columns=fullflexible,
  backgroundcolor=\color{gray!8},
  commentstyle=\color{matd48verde},
  keywordstyle=\color{matd48azul}\bfseries,
  language=R,
  extendedchars=true,
  literate=%
    {á}{{\'a}}1 {à}{{\`a}}1 {â}{{\^a}}1 {ã}{{\~a}}1
    {é}{{\'e}}1 {ê}{{\^e}}1 {í}{{\'i}}1 {ó}{{\'o}}1
    {ô}{{\^o}}1 {õ}{{\~o}}1 {ú}{{\'u}}1 {ç}{{\c c}}1
    {Á}{{\'A}}1 {À}{{\`A}}1 {Â}{{\^A}}1 {Ã}{{\~A}}1
    {É}{{\'E}}1 {Ê}{{\^E}}1 {Í}{{\'I}}1 {Ó}{{\'O}}1
    {Ô}{{\^O}}1 {Õ}{{\~O}}1 {Ú}{{\'U}}1 {Ç}{{\c C}}1
}

\hypersetup{colorlinks=true, linkcolor=matd48azul, urlcolor=matd48azul, citecolor=matd48azul}

\newcommand{\cabecalho}[2]{%
  \begin{center}
    {\Large\bfseries\color{matd48azul} #1}\\[0.3em]
    {\large #2}\\[0.3em]
    \rule{\linewidth}{0.6pt}
  \end{center}
  \vspace{0.5em}
}

\newenvironment{questao}[1]{\vspace{1em}\noindent\textbf{Questão #1.}\hspace{0.4em}}{\par}
\newenvironment{solucao}{\vspace{0.3em}\noindent\textit{Solução.}\hspace{0.4em}\color{black!85}}{\par\vspace{0.5em}}


\begin{document}

\cabecalho{Lista de Exercícios 12}{Fatoriais A×B, A×B×C e blocos completos (Capítulo 5 / Aula 12)}

\noindent \textit{Entrega: até a próxima aula. Justifique todas as respostas --- nestas listas, a
justificativa vale mais que a resposta final. O gabarito completo está em \texttt{Gabarito12.pdf}.}

\begin{questao}{1} \textbf{(Ciência de dados --- ANOVA de um fatorial A×B na mão)}

Uma equipe de engenharia mede o tempo de resposta médio (ms) de um endpoint sob dois algoritmos
de cache ($A$: LRU vs. LFU) cruzados com dois tamanhos de payload ($B$: pequeno vs. grande), com
$r=2$ execuções independentes por combinação:

\begin{center}
\begin{tabular}{c|cc}
 & \multicolumn{2}{c}{Payload ($B$)} \\
Algoritmo ($A$) & pequeno & grande \\ \hline
LRU & 10, 14 & 20, 24 \\
LFU & 16, 20 & 34, 38
\end{tabular}
\end{center}

\begin{enumerate}[label=(\alph*)]
  \item Calcule as quatro médias de célula, as duas médias marginais de $A$, as duas médias
    marginais de $B$ e a média geral.
  \item Usando as fórmulas $SQ_A = rb\sum_i(\bar y_{i\cdot\cdot}-\bar y_{\cdot\cdot\cdot})^2$,
    $SQ_B$ análogo, $SQ_{AB}=r\sum_{ij}(\bar y_{ij\cdot}-\bar y_{i\cdot\cdot}-\bar
    y_{\cdot j\cdot}+\bar y_{\cdot\cdot\cdot})^2$ e $SQ_E=\sum(y_{ijk}-\bar y_{ij\cdot})^2$,
    monte a tabela de ANOVA completa (gl, SQ, QM, F) para $A$, $B$, $A{\times}B$ e Erro.
  \item Ao nível $\alpha=0{,}05$, o valor crítico é $F_{1,4}=7{,}71$. Quais efeitos são
    significativos?
  \item Interprete o resultado: o algoritmo de cache e o tamanho do payload afetam o tempo de
    resposta de forma aditiva, ou o efeito de um depende do outro?
\end{enumerate}
\end{questao}

\begin{questao}{2} \textbf{(Agricultura --- interpretando uma ANOVA real)}

Um experimento cruza lâmina de irrigação (2 níveis) com dose de silício (4 níveis) sobre a altura
de plantas de pepino, com 3 repetições (24 parcelas). A ANOVA (ignorando blocos) resultou em:

\begin{center}
\begin{tabular}{lcccc}
\toprule
Fonte & gl & SQ & QM & $F$ \\ \midrule
Irrigação & 1 & 0,260 & 0,260 & 5,28 \\
Silício & 3 & 1,387 & 0,462 & 9,38 \\
Irrigação:Silício & 3 & 0,388 & 0,129 & 2,63 \\
Erro & 16 & 0,789 & 0,049 & \\ \bottomrule
\end{tabular}
\end{center}

\begin{enumerate}[label=(\alph*)]
  \item Ao nível de 5\%, quais fontes são estatisticamente significativas? (Use $F_{1,16} \approx
    4{,}49$ e $F_{3,16}\approx3{,}24$ como valores críticos aproximados.)
  \item O experimento também registrou 3 blocos de terreno. Ao incluir o bloco no modelo, o
    $QM$ do erro \emph{aumentou} ligeiramente, de 0,049 para 0,055. O que isso sugere sobre a
    escolha de blocar por terreno neste experimento específico?
  \item A blocagem, mesmo quando "não ajuda" como neste caso, ainda é uma decisão defensável de
    se tomar \emph{antes} de ver os dados? Justifique.
\end{enumerate}
\end{questao}

\begin{questao}{3} \textbf{(Dedução --- ortogonalidade no fatorial 2×2 e o estimador de
interação)}

Considere um fatorial $2\times2$ ($a=b=2$) com $r$ repetições por célula, fatores codificados em
efeito $\pm1$: $x_A \in \{-1,+1\}$ para o fator $A$, $x_B\in\{-1,+1\}$ para $B$, e a coluna de
interação $x_{AB} = x_A x_B$ (produto de Hadamard, Seção 5.4 do livro-texto).

\begin{enumerate}[label=(\alph*)]
  \item Mostre que $\mathbf{x}_A'\mathbf{x}_B = 0$ (ortogonalidade entre os efeitos principais).
    (Dica: liste as 4 combinações de $(x_A,x_B)$ e some os produtos $x_A x_B$ sobre as $4r$
    observações balanceadas.)
  \item Mostre, da mesma forma, que $\mathbf{x}_A'\mathbf{x}_{AB} = 0$ e
    $\mathbf{x}_B'\mathbf{x}_{AB} = 0$.
  \item Conclua que, neste desenho balanceado $2\times2$, $\mathbf{X}'\mathbf{X}$ (para
    $\mathbf{X}=[\mathbf{1}\,|\,\mathbf{x}_A\,|\,\mathbf{x}_B\,|\,\mathbf{x}_{AB}]$) é uma matriz
    diagonal.
  \item Usando a fórmula geral $\hat\beta_j = \mathbf{x}_j'\mathbf{Y}/\mathbf{x}_j'\mathbf{x}_j$
    válida quando $\mathbf{X}'\mathbf{X}$ é diagonal, mostre que o coeficiente de interação
    $\hat\beta_{AB}$ é proporcional ao contraste $\bar y_{11\cdot} - \bar y_{12\cdot} - \bar
    y_{21\cdot} + \bar y_{22\cdot}$, e identifique a constante de proporcionalidade em função de
    $r$.
\end{enumerate}
\end{questao}

\begin{questao}{4} \textbf{(Ciência de dados --- leitura causal de um teste A/B/n)}

Uma equipe de produto testa cor do botão ($A$: azul vs. laranja) cruzada com o dispositivo do
usuário ($B$: celular vs. desktop) sobre a taxa de conversão (\%). As médias de célula observadas
foram: (azul, celular) $=3{,}0$; (azul, desktop) $=5{,}0$; (laranja, celular) $=5{,}0$; (laranja,
desktop) $=5{,}0$.

\begin{enumerate}[label=(\alph*)]
  \item Calcule o efeito principal estimado da cor do botão (a diferença laranja $-$ azul,
    marginalizada sobre dispositivo).
  \item Calcule o efeito da cor \emph{separadamente} para cada dispositivo (celular e desktop).
    Os dois valores coincidem com o efeito principal do item (a)?
  \item Na linguagem da Seção 5.2 do livro-texto, o que esse padrão revela sobre a
    \textbf{heterogeneidade do efeito de tratamento}? Que decisão de produto (mostrar a mesma cor
    para todos vs. personalizar por dispositivo) esse padrão sugere?
\end{enumerate}
\end{questao}

\begin{questao}{5} \textbf{(Ciência de dados --- exercício computacional em R)}

O código abaixo ajusta o fatorial A×B do algoritmo de cache (Questão 1) usando \texttt{aov()} e
constrói a matriz de delineamento manualmente. Está incompleto.

\begin{lstlisting}
algoritmo <- factor(rep(c("LRU","LFU"), each = 4))
payload   <- factor(rep(rep(c("pequeno","grande"), each = 2), times = 2))
tempo     <- c(10, 14, 20, 24, 16, 20, 34, 38)
dados <- data.frame(algoritmo, payload, tempo)

# (a) complete para obter a matriz de delineamento com termo de interacao
X <- model.matrix(______________________, data = dados)

# (b) complete para confirmar que a coluna de interacao e o produto das principais
identical(X[, 4], ______________________)

# (c) ajuste o modelo e obtenha a tabela de ANOVA
modelo <- aov(______________________, data = dados)
summary(modelo)
\end{lstlisting}

\begin{enumerate}[label=(\alph*)]
  \item Complete as três linhas indicadas.
  \item A tabela de ANOVA produzida por \texttt{summary(modelo)} deve reproduzir os valores de
    $F$ que você calculou manualmente na Questão 1(b)--(c). Confirme (mentalmente ou rodando o
    código) que os graus de liberdade batem: quantos há para $A$, $B$, $A{:}B$ e Resíduos?
\end{enumerate}
\end{questao}

\end{document}
